\chapter{Complexity — Uniform Tableau Emitter Steps}
%%% generated by the blueprint pipeline from TCSlib.Complexity.CircuitComplexity.UniformTableauEmitterSteps
\section{Overview}

Declarations of the Lean module \texttt{TCSlib.Complexity.CircuitComplexity.UniformTableauEmitterSteps}.

\section{Declarations}

\begin{theorem}[Counter-program transition from a snapshot]
\lean{Complexity.UTab.goes_snap}
\label{Complexity.UTab.goes_snap}
\leanok
\uses{Complexity.CounterProg.Goes, Complexity.CounterProg.Goes.congr, Complexity.CounterProg.Goes.trans, Complexity.UTab.ClS.exit, Complexity.UTab.ClS.reg, Complexity.UTab.CpS.dst, Complexity.UTab.CpS.exit, Complexity.UTab.CpS.src, Complexity.UTab.ER, Complexity.UTab.ER.f, Complexity.UTab.ER.f_i, Complexity.UTab.ER.f_lb, Complexity.UTab.ER.set_lb, Complexity.UTab.Ki, Complexity.UTab.Str, Complexity.UTab.goes_clearR, Complexity.UTab.goes_copyR, Complexity.UTab.goes_instr_inv, Complexity.UTab.ibits, Complexity.UTab.prog, Complexity.UTab.snap_bits, Complexity.UTab.tpL, Complexity.snapIdx}
\difficulty{5}
Let $M$ be a finite-state Turing machine with Boolean symbols, let $x$ be a finite
Boolean string, and let $n,z,T,t$ be natural numbers with $z$ a finite Boolean string.
Suppose $t\le T$, and let $e$ be a record whose fields for the input length, hard-wired
length, step budget, current instruction, preceding instruction cell, and first
instruction of the current step satisfy respectively $e.nc=n$, $e.z=|z|$, $e.t=T$,
$e.i=\operatorname{snapIdx}(M,|z|+n,T,t)$, $e.q=0$ if $t=0$ and
$e.q=e.i-\operatorname{Str}(M,n,z,T)$ otherwise, and
$e.lb=t\operatorname{Str}(M,n,z,T)$; also suppose its scratch field is zero. For any
natural number $p_s$, starting at the snapshot-template label selected by whether $t=0$,
with register values from $e$ and input position $p_s$, the counter program reaches the
copy-loop label at phase zero if $t=0$, and the layout-processing label otherwise. It
leaves the input position unchanged, appends the bits encoding instruction $e.i$ to the
prior output, and has register values from the record obtained by updating $e.i$ to
$e.i+1$, $e.q$ to $0$ if $t=0$ and to $e.i+1-\operatorname{Str}(M,n,z,T)$ otherwise, and
$e.lb$ to $e.i+1$; this transition takes at most $\operatorname{Ki}(M)+2e.lb+7(e.i+1)+3$
steps.
\end{theorem}

\begin{definition}[Layer-loop starting record]
\lean{Complexity.UTab.lstart}
\label{Complexity.UTab.lstart}
\leanok
\uses{Complexity.UTab.ER, Complexity.UTab.Str}
Let $M$ be a finite-tape Turing machine, let $x$ be a finite Boolean word, and let
$n,z,T,t,\ell$ be nonnegative integers, with $z$ a finite Boolean word. Write
$S=S(M,n,z,T)$ for the associated stride. Define a record with fields
$(n,|z|,T,tS,q,tS,\ell,0,0,0,0,0)$, where $q=0$ if $t=0$ and $q=tS-S$ otherwise.
\end{definition}

\begin{definition}[Layer size bound]
\lean{Complexity.UTab.layerBound}
\label{Complexity.UTab.layerBound}
\leanok
\uses{Complexity.UTab.Ki, Complexity.UTab.Lmax, Complexity.UTab.Str}
Let $M$ be a finite-tape Turing machine, let $x$ be a finite Boolean word, and let $n,T$
be nonnegative integers and $z$ a finite Boolean word. Write $k(M),L(M),K(M)$ for the
associated machine quantities and $S=S(M,n,z,T)$ for the associated stride. Define the
layer bound to be $k(M)((2T+3)(L(M)+20))+(|z|+n+2)(L(M)+20)+K(M)+9((T+1)S)+6$.
\end{definition}

\begin{theorem}[Execution of a tableau layer]
\lean{Complexity.UTab.goes_layer}
\label{Complexity.UTab.goes_layer}
\leanok
\uses{Complexity.CounterProg.Goes, Complexity.CounterProg.Goes.congr, Complexity.CounterProg.Goes.trans, Complexity.CounterProg.goes_jz_pos, Complexity.UTab.ER.f, Complexity.UTab.Str, Complexity.UTab.goes_inputs, Complexity.UTab.goes_snap, Complexity.UTab.goes_tapes, Complexity.UTab.ibits, Complexity.UTab.layerBound, Complexity.UTab.lstart, Complexity.UTab.prog, Complexity.UTab.rT, Complexity.UTab.range'_add, Complexity.cfgStride, Complexity.snapIdx}
\difficulty{6}
Let $M$ be a finite-state Turing machine with Boolean tape symbols, let $x$ be a Boolean
input word, let $n$ be a natural number, let $z$ be a Boolean word, and let
$T,t,\ell,ps$ be natural numbers. Assume $1\le T$ and $t\le T$. Starting at the layer
label whose initial-step flag is true exactly when $t=0$, with the register values
specified by the layer-start record for $(n,z,T,t,\ell)$ and input position $ps$, the
counter program can execute in at most $B$ steps to the copy-loop label when $t=0$, and
to the layer label otherwise. Here $B=k(2T+3)(L+20)+( |z|+n+2)(L+20)+(L+3)+9(T+1)D+6$,
where $k$ is the number of work tapes, $L$ is the maximum template length, and $D$ is
the stride determined by $M,n,z,T$. The final register values are those of the
layer-start record for $(n,z,T,t+1,\ell)$, except that its $q$ register is set to $0$
when $t=0$ and otherwise to $(t+1)D-D$. The output is extended by the concatenation,
over all indices $i$ from $tD$ through but excluding $(t+1)D$, of the Boolean encoding
of the instruction for position $i$ in the tableau, and the execution property holds for
every output prefix.
\end{theorem}

\begin{theorem}[Successor start-record identity]
\lean{Complexity.UTab.lstart_succ}
\label{Complexity.UTab.lstart_succ}
\leanok
\uses{Complexity.UTab.Str, Complexity.UTab.lstart}
\difficulty{3}
Let $M$ be a finite-tape Turing machine with $k$ tapes, let $n,T,t,\text{lt}$ be natural
numbers, and let $z$ be a finite Boolean list. Set $S=k(2T+1)+|z|+n+3$. The start record
at step $t+1$ has $q$-coordinate $(t+1)S-S$, so replacing that coordinate by $(t+1)S-S$
leaves the record unchanged.
\end{theorem}

\begin{theorem}[Range partitioned into blocks]
\lean{Complexity.UTab.range_mul}
\label{Complexity.UTab.range_mul}
\leanok
\uses{Complexity.UTab.range'_eq_map}
\difficulty{4}
For natural numbers $a$ and $S$, the list $[0,1,\ldots,aS-1]$ is the concatenation, in
order of $t=0,1,\ldots,a-1$, of the lists $[tS,tS+1,\ldots,tS+S-1]$.
\end{theorem}

\begin{theorem}[Layered universal-tableau execution]
\lean{Complexity.UTab.goes_layers}
\label{Complexity.UTab.goes_layers}
\leanok
\uses{Complexity.CounterProg.Goes, Complexity.CounterProg.Goes.congr, Complexity.CounterProg.Goes.trans, Complexity.CounterProg.goes_loop, Complexity.UTab.CpS.dst, Complexity.UTab.CpS.src, Complexity.UTab.ER.f, Complexity.UTab.ER.update_f, Complexity.UTab.Hd.body, Complexity.UTab.Hd.exit, Complexity.UTab.Hd.reg, Complexity.UTab.Str, Complexity.UTab.goes_copyR, Complexity.UTab.goes_layer, Complexity.UTab.ibits, Complexity.UTab.layerBound, Complexity.UTab.lstart, Complexity.UTab.lstart_succ, Complexity.UTab.prog, Complexity.UTab.rLT, Complexity.UTab.range_mul, Complexity.UTab.tpL}
\difficulty{6}
Let $M$ be a Boolean multi-tape Turing machine, let $x$ be a finite Boolean input word,
and let $n,z,T,p$ be natural numbers with $1\le T$. The eleven-register counter program
constructed from $M$, run on $x$, has a run of at most
$(T+1)(\mathrm{layerBound}(n,z,T)+7T+10)$ steps from its layer-start label with the
register values specified by $\mathrm{lstart}(0,0)$ and input position $p$ to the label
for the final template, with register values specified by $\mathrm{lstart}(T+1,0)$ and
the same input position; the output appended during this run is the concatenation, over
indices $i=0,\ldots,(T+1)\mathrm{Str}(n,z,T)-1$, of the instruction-bit strings for
index $i$.
\end{theorem}
